\documentclass[a4paper]{article}
\usepackage{color}
\usepackage{url}
\usepackage[T2A]{fontenc}
\usepackage[utf8]{inputenc}
\usepackage{graphicx}
\usepackage{amsmath}
\usepackage{amssymb}

\usepackage[english,serbian]{babel}

\usepackage[unicode]{hyperref}
\hypersetup{colorlinks,citecolor=green,filecolor=green,linkcolor=blue,urlcolor=blue}

\begin{document}

\title{Povlačenje lanca uvis\\ \vskip 0.2in \small{Seminarski rad u okviru kursa\\ Osnove matematičkog modeliranja\\ Matematički fakultet} \vskip 5.0in}

\author{
Matija Đorđević\\
59/2022
\and
Filip Nedeljković\\
43/2022
}
\date{\today}
\maketitle
\newpage


\section{Opis problema}
Na zemlji leži sklupčan lanac velike dužine i težine T po dužnom metru. U početnom trenutku sa
nivoa zemlje uvis konstantna sila F deluje na lanac. Lanac počinje da se penje uvis.
Izvodimo matematički model kretanja vrha lanca. \\
Izračunavamoo maksimalnu visinu do koje će lanac stići i vreme potrebno da stigne do te visine, kao 
i vreme za koje će lanac ponovo pasti na zemlju. \\

\section{Fizicka formulacija}

Jedine sile koje deluju na lanac su konstantna sila F  koja vuče vrh lanca 
naviše i sila Zemljine teže (Fg). U posmatranom modelu nećemo razmarati
otpor vazduha i trenje podloge. Pretpostavljamo da je da je gravitaciono
ubrzanje g konstantno. Odbacujemo uticaj rotacije Zemlje na lanac.
Lanac je dovoljno dugačak tako da se ceo lanac nikada neće podići
sa Zemlje. \\
Na početku  se lanac nalazi na visini od 0m. Početna brzina kojom se kreće lanac je 0m/s. 
Ovaj model predstavlja sistem promenljive mase. 
Početna masa koja se podiže je 0kg. Kako se lanac kreće naviše 
tako se menja količina mase koju drži vrh lanca.\\
Da bi smo mogli opisati kretanje lanca koristimo Drugi Njutnov zakon.
Kako je masa promenljiva ne moze da se koristi Drugi njutnov zakon
u obliku F=ma, gde je F sila koja deluje na telo, m masa tela, a
a njegovo ubrzanje, nego ćemo silu koja deluje na telo izraziti 
preko promene impulsa kroz vreme, odnosno F=dp/dt, gde 
p predstavlja impuls tela, a t vreme.

\section{Formiranje modela}

\[
x\left( 0 \right)=0,\quad x'\left( 0 \right)=1,
\]
\[
F \text{-- konstantna sila uvis},\quad
T \text{-- težina po dužnom metru},\quad
g \text{-- ubrzanje gravitacije}
\]

\[
m\left(x\right) = \frac{T}{g}\,x
\]

\[
\frac{d}{dt}\left(mv\right)=F-Tx \text{- koristimo II Njutnov zakon preko impulsa}
\]

Promena impulsa je jedna 
razlici F i T*x gde x prdstvalja trenutnu visinu vrha lanca
Sređivanjem dobijamo


\[
\frac{T}{g}\frac{d}{dt}\left(x \, x'\right) = F - Tx
\]


\[
\frac{d}{dt}\left(x \, x'\right) = x'^{2} + xx'' 
\]

Iz prethodna dva koraka sledi:

\[
x'^{2} + xx'' = g\left(\frac{F}{T} - x\right)
\]

Transformišemo diferencijalnu jednačinu u lepši oblik pomoću

\[
x'' = v\frac{dv}{dx}  \quad \left(v = x' \quad v\text{ je brzina}\right) \quad \left(v = \frac{dx}{dt}\right) \quad \left(\frac{dv}{dt} = \frac{dvdx}{dxdt}\right) 
\]

\[
xv\frac{dv}{dx} + v^{2} = \frac{Fg}{T} - gx \quad \text{pomnožimo obe strane jednačine sa } 2x
\]

\[
2x^{2}v\frac{dv}{dx} + 2xv^{2} = \frac{2Fgx}{T} - 2gx^{2} 
\]

Kako je: 
\[
\frac{d}{dx}\left(x^{2}v^{2}\right) = 2xv^{2} + 2x^{2}v\frac{dv}{dx}
\]

Sledi da je:

\[
\frac{d}{dx}\left(x^{2}v^{2}\right) = \frac{2Fgx}{T} - 2gx^{2} \quad \left(\text{sada integralimo obe strane jednačine sa} \quad \int_0^x \right)
\]

\[
x^{2}v^{2} = \int_0^x \left(\frac{2Fgy}{T} - 2gy^{2}\right)dy
\]

\[
x^{2}v^{2} = \frac{Fgx^{2}}{T} - \frac{2gx^{3}}{3}  \quad \left(\text{podelimo izraz sa } x^{2}\right)
\]

\[
v^{2} = \frac{Fg}{T} - \frac{2gx}{3}
\]

Posmatraćemo slučaj kada lanac se diže uvis tako da uzimamo
pozitivno rešenje brzine

\[
v = \sqrt{\frac{Fg}{T} - \frac{2gx}{3}}
\]

\[
\frac{dx}{dt} = \sqrt{\frac{Fg}{T} - \frac{2gx}{3}}
\]

\[
dt = \frac{dx}{\sqrt{\frac{Fg}{T} - \frac{2gx}{3}}} \quad \left(\text{integralisemo od početnog trenutka 0 do nekog trenutka t} \quad \int_0^t \right)
\]

\section{Rezultati}




\end{document}




\end{document}
